\section{Automated Neural Segmentation}

The dominant approach to large-scale neural segmentation uses convolutional
neural networks trained to predict affinities or boundary probabilities, followed
by a graph-based merge step (watershed, multicut, or similar) to produce
instance labels. Flood-filling networks (FFN)~\cite{januszewski2018high} achieve
state-of-the-art results on EM data by iteratively extending segment predictions
from seeds, and are used as the baseline segmenter for the XPRESS challenge.
PyTorch Connectomics~\cite{lin2021pytorch} provides a flexible framework for
training and deploying multiple segmentation architectures.

\section{Proofreading Infrastructure}

The CAVE (Connectome Annotation Versioning Engine)~\cite{dorkenwald2022cave}
system provides infrastructure for storing, versioning, and updating connectomics
annotations at scale. CAVE implements a chunk-based merge graph (PyChunkedGraph)
that allows atomic merge and split operations with full provenance tracking.
Our pipeline is designed to produce outputs compatible with CAVE-style
infrastructure: candidate connections in CSV form map naturally to proposed
merge operations, and the Neuroglancer annotation export can be overlaid in
CAVE's Neuroglancer integration.

The MICrONS Consortium~\cite{microns2025} recently published a functional
connectome of mouse visual cortex ($\sim$200{,}000 neurons, $\sim$500 million
synapses), representing the current scale of ambition for automated connectomics.
The proofreading pipeline underlying MICrONS relied on CAVE and extensive human
review, highlighting the importance of automated tools that reduce the human
burden.

\section{Error Detection and Correction}

Zung et al.~\cite{zung2017} proposed a framework for detecting and correcting
segmentation errors using a classifier trained on local volume patches. Their
approach operates at the voxel level, re-examining the raw image near putative
error sites. Our work operates entirely in the post-segmentation graph domain,
which is faster (no re-examination of raw data) and format-agnostic (works on
any segmentation, not only those for which raw data is available).

Plaza et al.~\cite{plaza2018analyzing} provide a rigorous treatment of
segmentation quality metrics for connectomics, including the definitions of
split and merge errors, the rand index, and the variation of information used
in challenge evaluations. Their framework informed the precision/recall
accounting used in this work, particularly the handling of ambiguous (AMBIGUOUS)
outcomes, which are excluded from precision/recall denominators.

\section{Skeleton-Based Proofreading}

Skeleton-based proofreading approaches identify potential split errors by
finding pairs of skeletons whose endpoints are spatially close and geometrically
compatible. Early systems used simple distance thresholds; modern approaches
add learned classifiers. The key contribution of this work is the skeleton-node
(interior) indexing strategy, which extends the detectable split positions from
fragment endpoints to any point along the axon.

The XPRESS challenge~\cite{xpress2023} was specifically designed to benchmark
automated proofreading of myelinated axons in XNH data, providing standardized
ground-truth oracle pairs for training and evaluation. To our knowledge, this
work is among the first applications of a conservative, rule-based graph-level
pipeline to the XPRESS benchmark, achieving Recall = 0.9958 on the training volume
and Recall = 0.9879 on the held-out validation volume.

\section{Graph Construction Strategies}

Several works have explored variants of the fragment graph construction problem.
Contact-based graphs (edges between touching fragments) are simple but miss
genuine splits with small gaps. Proximity-based graphs (edges between fragments
within a distance threshold) are more general; the key design choice is what
set of points to index. Our work systematically compares endpoint-only vs.\
full skeleton-node indexing and quantifies the impact of each on oracle coverage,
finding that full skeleton-node indexing is essential for interior splits.

The long-range endpoint pass introduced in Experiment~11 is similar in spirit
to approaches used in EM proofreading that specifically target ``broken wire''
errors by searching for endpoint pairs at larger radii than the standard
neighborhood. Our contribution is the distance-conditioned scoring that
accompanies this pass, preventing the proximity feature from suppressing
genuine long-range candidates.
